% BHU PYQ -- Transcribed & Corrected
% Paper: ESMN111 / ESMJ11: Elements of Earth Sciences
% Exam: B.Sc. (Hons.) Semester I Examination, 2024-25 (NEP)
% Department: Department of Geology / Earth Science

\documentclass[11pt,a4paper]{article}
\usepackage[a4paper,top=2.5cm,bottom=2.5cm,left=2.5cm,right=2.5cm,headheight=14pt]{geometry}
\usepackage{amsmath,amssymb}
\usepackage[T1]{fontenc}
\usepackage[utf8]{inputenc}
\usepackage{lmodern,microtype,fancyhdr,enumitem,hyperref,lastpage,parskip}

\hypersetup{colorlinks=true,urlcolor=black,linkcolor=black,pdftitle={ESMN111: Elements of Earth Sciences -- BHU Sem I 2024-25 (NEP)}}
\pagestyle{fancy}\fancyhf{}
\fancyhead[L]{\small   \textit{Banaras Hindu University}}
\fancyhead[R]{\small   \textit{ESMN111: Elements of Earth Sciences}}
\fancyfoot[C]{\small Page  \thepage\ of \pageref{LastPage}}


\newlist{parts}{enumerate}{2}
\setlist[parts,1]{label=(\alph*),leftmargin=*,itemsep=4pt,topsep=2pt}
\setlist[parts,2]{label=(\roman*),leftmargin=*,itemsep=3pt,topsep=2pt}


\begin{document}


\begin{center}
  {\large\bfseries BANARAS HINDU UNIVERSITY}\\[2pt]
  {
\normalsize B.Sc. (Hons.) Semester I Examination 2024-25 (NEP)}\\[6pt]
  \rule{0.8\linewidth}{0.5pt}\\[6pt]
  {\large\bfseries DEPARTMENT OF GEOLOGY / EARTH SCIENCE}\\[4pt]
  {\large\bfseries Paper Code: ESMJ11 / ESMN111 --- Elements of Earth Sciences}\\[4pt]
  {
\normalsize Time: 2:30 Hours\hfill Maximum Marks: 70}
\end{center}

\rule{\linewidth}{0.4pt}\smallskip



\noindent   \textit{Instruction: Write your Roll No. at the top immediately on the receipt of this question paper.}\\[2pt]



\noindent   \textit{Note: Attempt five questions in all, selecting at least two questions each from Section--A and Section--B. Section--C is compulsory. All questions carry equal marks (14 marks each).}
\bigskip

\section*{SECTION -- A}




\noindent   \textbf{Question 1} \hfill [14 Marks]\\
Write descriptive notes about Nebular Hypothesis and Protoplanet Hypothesis.

\bigskip




\noindent   \textbf{Question 2} \hfill [14 Marks]\\
What is the concept of Uniformitarianism? Also write explanatory notes on the concepts of Neptunism, Vulcanism and Plutonism.

\bigskip




\noindent   \textbf{Question 3} \hfill [14 Marks]\\
Differentiate between:

\begin{parts}
  \item Continental crust and Oceanic crust
  \item Mud-slide and Debris-slide
  \item Relative dating and Absolute dating
  \item Strike joint and Dip joint
\end{parts}

\bigskip
\section*{SECTION -- B}




\noindent   \textbf{Question 4} \hfill [14 Marks]\\
What are the causes of volcanic eruptions? With neat diagrams, discuss the classification of volcanoes.

\bigskip




\noindent   \textbf{Question 5} \hfill [14 Marks]\\
What is a fold? With neat diagram describe different parts of a fold. Also describe the classification of folds based on:

\begin{parts}
  \item Upward and downward closing of folds
  \item Interlimb angle of folds
\end{parts}

\bigskip




\noindent   \textbf{Question 6} \hfill [14 Marks]\\
Write critical notes on any   \textbf{two} of the following:

\begin{parts}
  \item Propagation of P- and S-waves inside the Earth
  \item Principle of Superposition and its use in field geology
  \item Causes of earthquakes
  \item Fault classification based on the rake of the net slip
\end{parts}

\bigskip
\section*{SECTION -- C (Compulsory)}




\noindent   \textbf{Question 7} \hfill [14 Marks]\\
Describe the following in two or three sentences each:

\begin{parts}
  \item Mass of the Earth
  \item Repetti discontinuity
  \item NiFe
  \item Catastrophism
  \item Weathering
  \item Rayleigh waves
  \item Palaeozoic era
  \item Half-life of $^{14} \text{C}$
  \item Caldera
  \item True dip
  \item Antiformal syncline
  \item Throw of fault
  \item Outcrop
  \item Contour line
\end{parts}

\end{document}
